\section{方法}
\label{sec:method}

本节给出我们的核心贡献。我们先将第~\ref{sec:problem:objective} 节的在线选路优化形式化为 MDP;随后说明训练算法的选择;再介绍流感知的模型结构--它通过特征级线性调制(feature-wise linear modulation, FiLM)把每条流的 QoS 意图注入图神经网络的每一层消息传递(即节点反复聚合邻居信息、更新自身表示的过程);最后以伪代码给出完整训练流程。本节涉及的网络维度、层数与训练系数均以符号表示,其具体取值统一列于第~\ref{sec:exp:hyperparam} 节。

\subsection{问题形式化:MDP 建模}
\label{sec:method:mdp}

我们将第~\ref{sec:problem:objective} 节的在线选路优化形式化为马尔可夫决策过程 $(\mathcal{S},\mathcal{A},P,r)$。智能体(即流感知 FiLM-GNN 策略)在序列内逐流决策:第 $t$ 步为第 $t$ 条流 $\mathbf{f}_t$ 从其候选集 $\mathcal{P}_t$ 中选定一条路径。在每一步,策略据当前状态 $s_t$ 输出动作 $a_t$,环境执行选路、扣减带宽并返回奖励 $r_t$ 与下一状态 $s_{t+1}$。

\textbf{状态} $s_t\in\mathcal{S}$ 由三部分组成:
\begin{enumerate}
\item \emph{图快照}:节点特征(标记当前流的源宿 $o_t,v_t$)与链路特征 $[d_e,\,c_e,\,b_e^{\text{res}}]$(时延、容量、剩余带宽);
\item \emph{流意图向量} $m_t=[\phi_t,\,\tilde b_t]$,其中 $\tilde b_t$ 为归一化带宽需求;
\item \emph{候选路径} $\mathcal{P}_t$ 及其有效性掩码。
\end{enumerate}
观测仅包含当前流 $\mathbf{f}_t$,不含未来流 $\{\mathbf{f}_\tau\}_{\tau>t}$,与第~\ref{sec:problem:objective} 节的在线因果约束一致。

\textbf{动作} $a_t\in\mathcal{A}=\{0,1,\dots,K-1\}$ 为离散选择,从 $K$ 条候选路径中选定第 $a_t$ 条 $p_t=p_t^{a_t}$ 承载当前流。

\textbf{转移} 需区分两个层面。给定当前流与所选路径,\emph{资源记账是确定且已知的}:沿途链路的剩余带宽 $b_e^{\text{res}}$ 扣减 $\hat b_t$(式~\eqref{eq:rho},即实际分配带宽),指针前移。然而下一条到达的流 $\mathbf{f}_{t+1}$ 由需求分布 $\mathcal{D}$ 产生、在到达前对智能体不可见,且不受当前动作影响。因此完整状态的转移 $P(s_{t+1}\mid s_t,a_t)$ 并非确定性的(给定 $s_t,a_t$ 后 $s_{t+1}$ 不唯一)--其随机性来自下一条流的到达。这正是第~\ref{sec:problem:why_rl} 节所述"确定记账、未来不可见"的结构。

\textbf{奖励} 直接取自第~\ref{sec:problem:objective} 节的效用函数 $U(\mathbf{f}_t,p_t)$(式~\eqref{eq:utility})。平移一个常数 $-1$ 以收窄奖励范围、便于价值函数学习:
\begin{equation}
r_t = U(\mathbf{f}_t, p_t) - 1
    = -1 + U^{\text{bw}}_t + U^{\text{dl}}_t.
\label{eq:reward}
\end{equation}
奖励区间为 $[-1, 1]$:大象流($\phi_t\to 1$)的时延效用项趋于零,其奖励上界为 $0$($\rho_t=1$ 时取到);老鼠流($\phi_t\to 0$)在带宽被完全满足且时延极低时两项效用均可取到较大正值,使该类流的奖励可超过 $0$--因此在老鼠流占主导的测试分布下,均值奖励整体可为正。由于平移常数不改变最优策略,最大化累积奖励 $\sum_t r_t$ 与最大化累积效用(式~\eqref{eq:objective})等价。式~\eqref{eq:reward} 的结构意味着:智能体若要同时在老鼠流(要求带宽硬满足与低时延)与大象流(看重吞吐比率)上获得高奖励,其策略必须\emph{感知并区分每条流的 QoS 意图} $\phi_t$。这一要求正是下一节流感知表示的设计动机。

\subsection{训练算法的选择}
\label{sec:method:algo}

第~\ref{sec:problem:objective} 节论证了本问题是一个在需求分布上的在线随机优化,需要非短视策略。我们选择近端策略优化(Proximal Policy Optimization, PPO)作为训练算法:其随机策略以 Categorical 分布天然参数化从 $K$ 条候选路径中的离散选择,裁剪式目标使更新稳定(契合两类流奖励尺度差异较大的情形);同策略采集与在线部署设定一致,配合 GAE 实现非短视所需的长程信用分配;且作为无模型方法,无需对未知到达分布 $\mathcal{D}$ 导致的转移随机性显式建模。

\subsection{流感知 FiLM-GNN 编码器}
\label{sec:method:encoder}

编码器把状态中的图快照映射为节点嵌入（即每个节点的向量表示）。核心思想是：消息传递的每一层都应当"知道当前在为哪条流选路"，从而对不同 QoS 意图的流生成不同的拓扑表示。我们用 FiLM~\cite{Perez2018_FiLM} 实现这一条件化。记链路特征向量 $f_e = [d_e,\, c_e,\, b_e^{\text{res}}] \in \mathbb{R}^3$，即时延、容量与剩余带宽的拼接。

\begin{figure}[t]
\centering
\includegraphics[width=\linewidth]{FiLM GNN Reaction to Env.pdf}
\caption{FiLM-GNN 整体架构。左侧为流感知 FiLM-GNN 编码器：每层以当前流的 QoS 意图向量 $m_t = [\phi_t, \tilde b_t]$ 生成 FiLM 参数 $(\gamma^{(\ell)}, \beta^{(\ell)})$，经残差调制 $x^{(\ell)} = (1+\gamma^{(\ell)})\odot\tilde x^{(\ell)}+\beta^{(\ell)}+x^{(\ell-1)}$ 重塑节点嵌入，使同一拓扑对不同意图的流产生差异化表示。右侧为策略--价值双头结构：Actor Head 通过路径池化器（PathPooler，LSTM $+$ LayerNorm）与 $K$ 路打分器（KPathScorer）输出 Categorical 动作分布；Critic Head 以注意力池化聚合节点嵌入、经 MLP 回归状态价值 $V_\psi(s_t)$。右下方为 MDP 闭环与 PPO 训练流程。}
\label{fig:film_gnn_arch}
\end{figure}

设隐藏维为 $h$。首先用一个多层感知机 $\mathrm{MLP}_{\text{lift}}$ 将节点原始特征升维到 $h$ 维:$x^{(0)}=\mathrm{MLP}_{\text{lift}}(x)$。随后堆叠 $L$ 层图注意力块。第 $\ell$ 层接收当前流的意图向量 $m_t$,经两个独立的 MLP 生成逐通道(对 $h$ 维隐藏向量逐维)的缩放与平移参数:
\begin{align}
\gamma^{(\ell)} &= \mathrm{MLP}^{(\ell)}_{\gamma}(m_t), \quad \gamma^{(\ell)}\in\mathbb{R}^{h},
\label{eq:film_gamma}\\
\beta^{(\ell)} &= \mathrm{MLP}^{(\ell)}_{\beta}(m_t), \quad \beta^{(\ell)}\in\mathbb{R}^{h},
\label{eq:film_beta}
\end{align}
并对消息传递的输出施加 FiLM 调制:
\begin{align}
\tilde x^{(\ell)} &= \mathrm{ELU}\!\big(\mathrm{LN}(\mathrm{Conv}^{(\ell)}(x^{(\ell-1)}, E, f))\big),
\label{eq:conv}\\
x^{(\ell)} &= \big(1+\gamma^{(\ell)}\big)\odot \tilde x^{(\ell)} + \beta^{(\ell)} \;+\; x^{(\ell-1)},
\label{eq:film}
\end{align}
其中 $\mathrm{Conv}^{(\ell)}$ 为图注意力算子(默认 TransformerConv~\cite{Shi2021_UniMP} 风格的点积注意力,以链路特征 $f$ 作为边属性),$\mathrm{LN}$ 为层归一化,$\odot$ 为逐元素乘,最后一项为残差连接(把该层输入直接加到输出,缓解深层网络的训练困难)。式~\eqref{eq:film} 是本方法的关键:意图向量 $m_t$ 在\emph{每一层}都以仿射变换 $(1+\gamma)\odot\tilde x+\beta$ 重塑节点表示,使得同一张物理拓扑在面对老鼠流与大象流时被编码成不同的嵌入。这与 GNN-FiLM~\cite{Brockschmidt2020_GNNFiLM} 的根本区别在于:后者用\emph{目标节点表示}计算 $\gamma,\beta$ 以调制聚合,而我们用\emph{逐流的 QoS 意图}计算 $\gamma,\beta$,使消息传递成为流感知(flow-aware)的。

\subsection{路径池化、打分与价值网络}
\label{sec:method:scorer}

\textbf{路径池化。} 节点嵌入 $x^{(L)}$ 经一个路径池化器(PathPooler)转换为候选路径表示。对第 $k$ 条候选路径,沿其节点序列取出嵌入构成变长序列,用一个 LSTM 编码并取末隐状态、再经层归一化,得到路径嵌入 $z_k\in\mathbb{R}^{h}$。这样 $K$ 条候选路径被池化为 $\{z_k\}_{k=1}^{K}$。

\textbf{打分与策略。} 打分器 $\mathrm{KPathScorer}$ 把每条路径嵌入与当前流意图拼接 $[z_k; m_t]$,先经特征 MLP(ELU $+$ LayerNorm)再经打分头 MLP 输出标量分值 $\ell_k$(logit,即未归一化的打分)。$K$ 个分值构成一个 Categorical 分布,策略据此采样动作:
\begin{equation}
\pi_\theta(a_t\!=\!k\mid s_t) = \mathrm{softmax}_k(\ell_1,\dots,\ell_K).
\label{eq:policy}
\end{equation}
注意,若意图向量 $m_t$ 仅在此处(末端打分器)一次性进入,此时路径已被池化为固定向量,意图信息无法再反向影响拓扑表示,造成不可逆的信息损失;这正是本方法选择在\emph{每一层}注入意图、而非仅在末端拼接的动机。

\textbf{价值网络。} 价值网络(StateCritic)使用一份独立的编码器,对其输出的节点嵌入用一个以意图为条件的注意力池化(按学习到的权重对各节点嵌入加权求和,汇聚为单个图级向量)得到图级嵌入,再经 MLP 回归为标量状态价值 $V_\psi(s_t)$。策略网络与价值网络不共享编码器参数。

\subsection{训练流程}
\label{sec:method:train}

训练采用 PPO,优势用广义优势估计(GAE)计算,损失由裁剪策略目标 $\mathcal{L}^{\text{clip}}_\theta$、价值损失 $c_v\mathcal{L}^{\text{value}}_\psi$ 之和减去熵奖励项 $c_e\mathcal{H}[\pi_\theta]$ 构成:
\begin{equation}
\mathcal{L}(\theta,\psi) = \mathcal{L}^{\text{clip}}_\theta + c_v\,\mathcal{L}^{\text{value}}_\psi - c_e\,\mathcal{H}[\pi_\theta],
\label{eq:loss}
\end{equation}
其中 $\mathcal{L}^{\text{clip}}_\theta$ 为裁剪系数 $\epsilon$ 控制的策略目标,$c_v,c_e$ 分别为价值与熵系数,优势在 minibatch 内标准化。完整流程见算法~\ref{alg:ppo}。

\begin{algorithm}[t]
\caption{流感知 FiLM-GNN 的 PPO 训练}
\label{alg:ppo}
\SetKwInOut{Input}{输入}
\SetKwInOut{Output}{输出}
\Input{网络 $G$;需求分布 $\mathcal{D}$;候选数 $K$;序列长度 $W$;总帧数 $T$;每批帧数 $B$;PPO 轮数 $E$;学习率 $\eta$}
\Output{策略参数 $\theta$、价值参数 $\psi$}
随机初始化 $\theta,\psi$\;
\While{已采集帧数 $<T$}{
  \tcp{采集(同策略)}
  用当前策略 $\pi_\theta$ 与环境交互采集 $B$ 帧;每步:编码图$\to$池化 $K$ 路径$\to$打分$\to$按 $\pi_\theta$ 采样动作$\to$执行选路、扣减带宽、记奖励 $r_t$(式~\eqref{eq:reward})\;
  \tcp{优势估计}
  用 $V_\psi$ 与 GAE 计算优势 $\hat A_t$ 与回报目标 $\hat R_t$\;
  \tcp{优化}
  \For{$E$ 个 PPO epoch}{
    \ForEach{minibatch}{
      在 minibatch 内标准化 $\hat A_t$\;
      计算损失 $\mathcal{L}(\theta,\psi)$(式~\eqref{eq:loss})\;
      梯度裁剪后用 Adam(学习率 $\eta$)更新 $\theta,\psi$\;
    }
  }
  学习率按余弦退火更新\;
  \tcp{评估与保存}
  以确定性策略跑若干 rollout,记录平均奖励;保存最优检查点\;
}
\end{algorithm}
